% ISMRM abstract: MGLPDSNet, a multigrid extension of LPDSNet.
% Structure follows the ISMRM submission form: Impact, Synopsis (Motivation / Goal(s) / Approach /
% Results), Abstract body (Introduction ... Conclusion), Acknowledgements, References, Figures.
% The form takes plain text pasted section by section; this file is for drafting, word counts and
% previewing the figures, not for upload.
%
% \todo{...} marks every number or decision still open. Every setting in Methods is read from the
% generated configs (scripts/make_mg_recon_configs.py; the exp16 cells *vtxAnestU and their baselines).
% Limits (ISMRM, check the current call): Impact <= 40 words, Synopsis <= 100 words,
% body <= 750 words, <= 5 figures. Word counts: texcount main.tex.

\documentclass[10pt]{article}
\usepackage[margin=2cm]{geometry}
\usepackage[T1]{fontenc}
\usepackage{lmodern}
\usepackage{amsmath,amssymb,bm}
\usepackage{graphicx}
\usepackage{booktabs}
\usepackage{xcolor}
\usepackage{tikz}
\usetikzlibrary{arrows.meta,calc,positioning}
\usepackage{pgfplots}
\usepgfplotslibrary{groupplots}
\pgfplotsset{compat=1.16}
\usepackage{caption}
\captionsetup{font=small,labelfont=bf}
\usepackage[numbers,super,sort&compress]{natbib}

\newcommand{\todo}[1]{\textcolor{red}{[#1]}}
\newcommand{\figdir}{figures}
% a framed placeholder until a figure file exists
\newcommand{\figorbox}[2]{\IfFileExists{#1}{\includegraphics[width=\linewidth]{#1}}%
  {\fbox{\parbox[c][4cm][c]{0.95\linewidth}{\centering\small\textcolor{red}{#2}}}}}

\begin{document}

% ------------------------------------------------------------------------------------------
\begin{center}
{\Large\bfseries Multigrid Learned Primal-Dual Splitting\\
for Artifact-Free Accelerated MRI Reconstruction}\\[6pt]
Eric Eng, \todo{co-authors}, Yao Wang, Li Feng
\end{center}

% ------------------------------------------------------------------------------------------
\section*{Impact}
% <= 40 words
Compact unrolled networks for accelerated MRI occasionally leave residual aliasing artifacts that
undermine trust in their reconstructions. A multigrid extension suppresses these artifacts at
matched size and runtime, making algorithm-derived reconstruction networks more reliable.

\section*{Synopsis}
\textbf{Motivation:} Directly parameterized unrolled networks such as LPDSNet achieve
state-of-the-art accelerated MRI reconstruction with few parameters, but can leave residual
aliasing artifacts.\\
\textbf{Goal(s):} Suppress these artifacts while keeping LPDSNet's compact, algorithm-derived
architecture.\\
\textbf{Approach:} We embed LPDSNet's primal-dual iterations in a multigrid V-cycle derived from
MGProx. Coarse grids solve rediscretized versions of the same reconstruction problem, so spatially
extended aliasing is corrected cheaply; the network is trained end-to-end.\\
\textbf{Results:} On fastMRI knee and brain at R=8 and 12, MGLPDSNet \todo{removes the aliasing
artifacts} of LPDSNet at \todo{matched parameters and runtime}, while \todo{matching/exceeding}
E2E-VarNet.

% ------------------------------------------------------------------------------------------
\section*{Introduction}
Unrolled networks reconstruct accelerated MRI by alternating data consistency with a learned
prior. Directly parameterized networks such as CDLNet\cite{ref_cdlnet} and
LPDSNet\cite{ref_lpdsnet} go further: every layer is one iteration of a classical algorithm, here
primal-dual splitting, so the network is compact, noise-adaptive and interpretable by
construction, while matching black-box networks such as E2E-VarNet\cite{ref_varnet}. In some
reconstructions, however, LPDSNet leaves residual aliasing (Figure~\ref{fig:alias}). Undersampling
aliasing is spatially global, folding energy across the field of view, whereas each LPDS iteration
applies only local filters at full resolution, so a fixed budget of iterations cannot move
information far enough to undo it. Multigrid methods remove exactly this kind of spatially
extended error by correcting it on coarse grids\cite{ref_brandt}, and MGProx\cite{ref_mgprox}
extends them to the nonsmooth problems that LPDSNet solves. We propose MGLPDSNet, which embeds
LPDSNet's iterations in a multigrid V-cycle.

\section*{Methods}
\textbf{LPDSNet:} Given undersampled multicoil k-space $\bm y=\bm E\bm x+\bm n$, with
$\bm E=\bm M\bm F\bm S$ (sampling mask, Fourier transform, coil sensitivities), each LPDS layer takes
a primal gradient step on the image $\bm x$ and a dual step on latent subband coefficients $\bm z$
(Figure~\ref{fig:arch}): learned analysis and synthesis convolutions $\bm A^H$, $\bm B$ replace the
sparsifying transform, and the dual proximal map is a subband-wise clipping $\Pi_{\bm\lambda}$.
Step sizes and thresholds are affine in the noise level $\sigma$\cite{ref_cdlnet,ref_lpdsnet}.

\textbf{MGLPDSNet:} Each iteration is a V-cycle over three grid levels ($N$, $N/2$, $N/4$;
Figure~\ref{fig:arch}). LPDS layers at a level are followed by an FAS-Correction, which restricts
the state to the next coarser grid and tilts that level's problem so it reproduces the fine
level's residual: the full approximation scheme\cite{ref_brandt}, applied to both the primal and
the dual residual. Each coarse level solves the same SENSE problem posed on its own grid. Its
result returns through a Coarse Correction, prolonging the coarse update with a learned step,
before further LPDS layers. MGLPDSNet unrolls \todo{6} V-cycles with \todo{4/4/6} layers per level
and $M=\todo{81}$ subbands at every level, all weights untied.

\textbf{Experiments:} We trained on multicoil fastMRI\cite{ref_fastmri} knee (proton density) and
brain (T2-weighted) data, using measured k-space with added white Gaussian noise of
$\sigma\in[0.01,0.02]$ and 1-D random Cartesian undersampling at R=8 and 12 with a 4\% fully
sampled centre. Coil sensitivities were estimated with ESPIRiT from the noisy calibration region;
targets were SENSE combinations of the fully sampled data. All networks were trained identically
(magnitude $\ell_1$+$\ell_2$ loss within the anatomy mask, Adam, learning rate $5\times10^{-4}$
with cosine decay, $3\times10^5$ steps). Baselines were LPDSNet\cite{ref_lpdsnet} (30 layers,
169 subbands) and E2E-VarNet\cite{ref_varnet} given the same coil sensitivities
\todo{, and RAM\cite{ref_ram} trained under the same protocol}. PSNR, SSIM and NRMSE were computed
within the anatomy mask on 815 brain and 100 knee test volumes at $\sigma=\todo{0.015}$, with paired Wilcoxon
signed-rank tests (Holm-corrected); network time was measured on one NVIDIA L40S.

\section*{Results}
Figure~\ref{fig:recon} shows representative reconstructions for each anatomy and acceleration.
\todo{One or two sentences: where MGLPDSNet is visibly better, and what the zoomed regions show.}

Figure~\ref{fig:alias} shows LPDSNet's residual aliasing: \todo{describe, e.g.\ a ghost of the
anatomy shifted by FOV/R, clearly visible in the error map}. MGLPDSNet reconstructs the same slices
without it. \todo{How often it occurs: found in N of M test slices by ...}

Across 815 brain and 100 knee test volumes (Figure~\ref{fig:bars}), MGLPDSNet improved on both
LPDSNet and E2E-VarNet in every metric, anatomy and acceleration ($p<0.001$), by
0.28--0.58\,dB PSNR over LPDSNet and 0.28--0.57\,dB over E2E-VarNet. It does so with 1.73M parameters, $17\times$ fewer than E2E-VarNet (29.5M), at an inference time of
50\,ms per slice, against 39\,ms for LPDSNet and 48\,ms for E2E-VarNet.

The ablations (Figure~\ref{fig:ablation}) show \todo{how quality scales with the number of
subbands $M$ and of V-cycles, and the single-resolution LPDSNet row as the no-multigrid reference}.

\section*{Discussion}
The coarse levels give every V-cycle a path along which information travels across the image in
few operations, which is what undoing global aliasing requires; local full-resolution iterations
lack it. Because each coarse level solves the same reconstruction problem, the network keeps the
structure of its source algorithm rather than becoming a generic multiscale CNN.
\todo{Limitations: retrospective undersampling, Cartesian sampling, a single noise model.}

\section*{Conclusion}
Embedding LPDSNet in a multigrid V-cycle suppresses its residual aliasing artifacts at
\todo{matched size and runtime}, while \todo{keeping} state-of-the-art reconstruction quality.

\section*{Acknowledgements}
\todo{Funding; NYU computing resources.}

{\small
\begin{thebibliography}{9}
\bibitem{ref_cdlnet} Janjušević N, Khalilian-Gourtani A, Wang Y. CDLNet: noise-adaptive
  convolutional dictionary learning network for blind denoising and demosaicing. IEEE Open J
  Signal Process. 2022;3:196--211. \todo{verify pages}
\bibitem{ref_lpdsnet} Janjušević N, Khalilian-Gourtani A, Wang Y, Feng L. Learned primal dual
  splitting for self-supervised noise-adaptive MRI reconstruction. arXiv:2504.15390. 2025.
  \todo{cite the published version}
\bibitem{ref_varnet} Sriram A, Zbontar J, Murrell T, et al. End-to-end variational networks for
  accelerated MRI reconstruction. MICCAI, LNCS 12262:64--73. 2020.
\bibitem{ref_brandt} Brandt A. Multi-level adaptive solutions to boundary-value problems. Math
  Comput. 1977;31(138):333--390.
\bibitem{ref_mgprox} Ang A, De Sterck H, Vavasis S. MGProx: a nonsmooth multigrid proximal gradient
  method with adaptive restriction for strongly convex optimization. SIAM J Optim. 2024.
  \todo{verify volume/pages}
\bibitem{ref_fastmri} Zbontar J, Knoll F, Sriram A, et al. fastMRI: an open dataset and benchmarks
  for accelerated MRI. arXiv:1811.08839. 2018.
\bibitem{ref_ram} Terris M, Hurault S, Song M, Tachella J. Reconstruct Anything Model: a lightweight
  foundation model for computational imaging. 2025. \todo{verify; drop if RAM is not reported}
\end{thebibliography}}

% ------------------------------------------------------------------------------------------
\clearpage
\section*{Figures}

\begin{figure}[h]
\centering
\resizebox{\linewidth}{!}{\input{\figdir/fig1_arch.tikz}}
\caption{MGLPDSNet. Top left: \todo{6} untied V-cycles, each receiving the preprocessed
measurement $\tilde{\bm y}=\bm E^H(\bm y-\mu\bm E\bm 1)$, $\bm E$ and $\sigma$; the first layer
starts from $\bm x=\tilde{\bm y}$, $\bm z^{(0)}=\Pi_{\bm\lambda}(\bm A^H\tilde{\bm y})$. Top right:
one V-cycle over three grid levels; each level solves its own problem
$(\tilde{\bm y}_l,\bm E_l,\sigma_l)$, tilted by the FAS correction $\bm\pi^{(l)}$. Middle: an LPDS
Layer, with $\tilde{\bm x}=\bm x^{+}+\theta(\bm x^{+}-\bm x)$; $\eta,\theta,\bm\lambda$ are affine in
$\sigma$; dashed inputs (the FAS terms) exist on coarse levels only. The Coarse Correction acts on
$\bm u\in\{\bm x,\bm z\}$ with $\hat{\bm u}_c$ the coarse level's output and $\bm u_c$ its
restricted input. Bottom: the FAS-Correction, in the order its operations run; it carries its own
learned copies of $\bm A^H$, $\bm B$ and $\bm\lambda$.}
\label{fig:arch}
\end{figure}

\begin{figure}[h]
\centering
\figorbox{\figdir/fig2_recon.pdf}{Fig.\ 2: reconstructions, rows knee R8 / knee R12 / brain R8 /
brain R12; columns LPDSNet, E2E-VarNet, \todo{RAM}, MGLPDSNet, reference; zoom and PSNR/SSIM per
panel; slices chosen by rule, not by eye.}
\caption{\todo{Representative reconstructions ...}}
\label{fig:recon}
\end{figure}

\begin{figure}[h]
\centering
\includegraphics[width=\linewidth]{\figdir/fig3_alias.png}
\caption{Residual aliasing in LPDSNet. Columns: zero-filled, E2E-VarNet, LPDSNet, MGLPDSNet and
the ground truth, with PSNR per slice. Below each image: the boxed region and its absolute error.
LPDSNet leaves thin line artifacts that are absent from the ground truth (brain, $R{=}8$: a bright
streak; brain, $R{=}12$: a dark line; knee, $R{=}12$: an oblique streak), visible as bright lines in
its error maps; MGLPDSNet and E2E-VarNet reconstruct the same regions without them.}
\label{fig:alias}
\end{figure}

\begin{figure}[h]
\centering
% main_table.tex: copy of figures/out/main_table.tex (rerun figures/table_main.py)
{\small\input{\figdir/main_table.tex}}\\[10pt]
\input{\figdir/fig4_timing.tikz}
\caption{Quantitative comparison. Top: PSNR, SSIM and NRMSE within the anatomy mask, mean
$\pm$ s.e.m.\ over test volumes (brain $n{=}815$, knee $n{=}100$); bold: best per experiment;
$^{*}$: differs from MGLPDSNet (paired Wilcoxon signed-rank over volumes, Holm-adjusted,
$p<0.05$). Bottom: network time per $640\times320$, 20-coil slice on one NVIDIA L40S (median of
30 runs), forward and backward pass. Parameters: E2E-VarNet 29.5M, LPDSNet 1.0M, MGLPDSNet
($M{=}81$) 1.7M.}
\label{fig:bars}
\end{figure}

\begin{figure}[h]
\centering
\input{\figdir/fig5_ablation_tradeoff.tikz}\\[10pt]
{\small\input{\figdir/ablation_table.tex}}
\caption{Ablations, knee R=8: time against quality. Top: test PSNR against inference time
($640\times320$, 20 coils, NVIDIA L40S) across the width sweep (subbands $M$ at 6 V-cycles, left)
and the depth sweep (V-cycles at $M{=}81$, right), with LPDSNet and E2E-VarNet for reference.
Bottom: the same networks in full, mean $\pm$ s.e.m.\ over $n{=}100$ test volumes; bold: best
in its sweep; $^{*}$: differs from $M{=}81$, 6 V-cycles (paired Wilcoxon signed-rank over volumes,
Holm-adjusted, $p<0.05$). \todo{4 V-cycles: still training}}
\label{fig:ablation}
\end{figure}

\end{document}
